%\documentclass[11pt]{article}

\documentclass[acmsmall,screen,nonacm]{acmart}
% 'acmsmall' = Guaranteed single column body text

% --- PACKAGES ---
\usepackage{filecontents}
\usepackage{graphicx}
\usepackage{hyperref}
\usepackage{geometry}
\usepackage{array}
\usepackage{fancyhdr}
\usepackage[margin=1in]{geometry}
\usepackage{parskip}
\usepackage{setspace}
\usepackage[hidelinks]{hyperref}

\usepackage{tabularx}
\usepackage{amssymb} % for \checkmark

% --- PACKAGES ---
\usepackage{filecontents}
\usepackage{graphicx}
\usepackage{hyperref}
\usepackage{geometry}
\usepackage{array}
\usepackage{fancyhdr}


\setstretch{1.08}
% --- DISABLE DEFAULT TITLE PROCESSING ---

\makeatletter
\makeatother

% Clear default ACM headers
\pagestyle{fancy}
\fancyhf{} 

\begin{document}


\begin{center}
    % TITLE
    {\Huge \textbf{Project Literature Survey :: Real-Time Wildfire Resilience Digital Twin}} \\
    \vspace{0.5em}
    {\Large \textbf{Group Name: N-Dimentional}} \\
    \vspace{2em}

    % AUTHORS TABLE
    % \hspace*{-2cm} pulls the table left so 4 columns fit easily
    \hspace*{-1.5cm} 
    \begin{tabular}{ >{\centering\arraybackslash}p{3.8cm} >{\centering\arraybackslash}p{3.8cm} >{\centering\arraybackslash}p{3.8cm} >{\centering\arraybackslash}p{3.8cm} }
        
        % ROW 1: Names
        \textbf{Prince Savsaviya}\textsuperscript{*} & \textbf{Viswanadh R. Challa}\textsuperscript{*} & \textbf{Anish Ashok Kale}\textsuperscript{*} & \textbf{Ali Rezayi Nejad}\textsuperscript{*} \\
        
        % ROW 2: IDs
        \small ID: 862548420 & \small ID: 862548873 & \small ID: 862547289 & \small ID: 862392323 \\
        
        % ROW 3: University
        \small UC Riverside & \small UC Riverside & \small UC Riverside & \small UC Riverside \\
        
        % ROW 4: Emails
        \small \href{mailto:psavs001@ucr.edu}{psavs001@ucr.edu} & 
        \small \href{mailto:vchal008@ucr.edu}{vchal008@ucr.edu} & 
        \small \href{mailto:akale014@ucr.edu}{akale014@ucr.edu} & 
        \small \href{mailto:areza018@ucr.edu}{areza018@ucr.edu} \\
        
    \end{tabular}
    
    \vspace{1em}
    % FOOTNOTE
    {\footnotesize \textsuperscript{*}Department of Computer Science, University of California, Riverside}
    \vspace{2em}
\end{center}

\date{February 19, 2026}


\section*{Abstract}
This literature survey supports a team project that builds a \emph{wildfire digital twin}: a continuously
updated virtual representation of wildfire conditions that ingests heterogeneous streams (satellite active-fire
detections, weather forecasts, fuel/topography layers, and optional ground sensors), performs real-time
spatiotemporal analytics, and produces actionable outputs such as risk indices, predicted spread envelopes,
and alerting dashboards.

We categorize prior work into: (1) digital-twin and dynamic data-driven foundations, (2) wildfire spread and
fire--atmosphere simulation models, (3) remote sensing fire detection products, (4) data assimilation and
real-time model updating, (5) cyberinfrastructure and decision-support platforms for wildfire, and (6) big-data
streaming and geospatial processing frameworks. Across these categories we identify a gap in end-to-end
reproducible pipelines that combine streaming ingestion, scalable geospatial joins, and model updating in a
way that is deployable by small teams. The proposed project targets that gap with a lightweight, open-source
architecture.

\section{Project context and key terms}
A \textbf{digital twin (DT)} is a digital representation of a physical system that stays synchronized with the
real system through data, enabling monitoring, forecasting, and decision support. Early DT definitions
emphasize a tight loop between the physical asset, its virtual counterpart, and the data connections between
them \cite{grieves2014digitaltwin}.

A wildfire DT is naturally a \textbf{dynamic data-driven application system (DDDAS)}, where streaming
observations are assimilated into running simulations and, optionally, simulations can influence where/what to
measure next (e.g., camera tasking). Darema formalized DDDAS as a paradigm that injects online data into an
executing model and can steer measurements in return \cite{darema2004dddas}.

\textbf{Wildfire Digital Twin objective for this project.} We aim to build an end-to-end pipeline that
(i) ingests near-real-time fire detections and environmental context, (ii) computes spatiotemporal risk layers
and simple spread forecasts, and (iii) serves results via an interactive map/dashboard and alerts. The remainder
of this survey groups related work into categories that directly map to these requirements.

\section{Literature selection scope}
We prioritized peer-reviewed or widely adopted technical reports that are frequently referenced in wildfire
science, remote sensing, or scalable data systems. Selection criteria: (a) practical influence on operational
wildfire modeling or monitoring, (b) relevance to real-time or data-driven updating, and (c) architectural
relevance to big-data streaming and geospatial analytics. The resulting set includes 20 works.

\section{Related work by category}

\subsection{Digital twin and dynamic data-driven foundations}
This category defines what it means to be a digital twin (continuous synchronization, actionable state
estimation) and why wildfire DTs require dynamic data-driven model updating.

\textbf{Digital Twin baseline.} Grieves describes the DT as a virtual representation tightly linked to the
physical system through data connections, enabling monitoring and prediction across the lifecycle
\cite{grieves2014digitaltwin}. For wildfire, the analogy is the evolving physical fire and a coupled virtual model
representing spread and risk.

\textbf{DDDAS paradigm.} Darema introduces DDDAS as applications that incorporate online data into executing
simulations and can dynamically steer measurement processes \cite{darema2004dddas}. Wildfire DTs are a direct
instantiation: ingest satellite detections and weather to update model state, and potentially request
higher-resolution imagery where uncertainty is high.

\textbf{Wildfire DT concept and roadmaps.} Halem et al.\ discuss a dynamic data-driven wildfire digital twin vision
emphasizing real-time data fusion, rapid prediction, and uncertainty-aware decision support \cite{halem2024wdt}.
Huang et al.\ summarize progress and challenges in wildfire digital twins (data latency, model fidelity,
assimilation, and compute constraints) and propose an integrated wildfire DT framing detection, simulation, and
prediction as a unified system \cite{huang2024wdtprogress}.

\textbf{Relevance to our project.} These works motivate our core loop:
(data ingestion) $\rightarrow$ (state estimation / risk layers) $\rightarrow$ (forecast) $\rightarrow$ (decision support),
and justify treating the system as a continuously updating DT rather than a one-off offline analysis.

\subsection{Wildfire spread and fire--atmosphere simulation models}
Wildfire DTs often rely on a simulation core to extrapolate from current state to near-future outcomes. This
category ranges from semi-empirical surface-spread models to coupled fire--atmosphere models.

\textbf{Surface fire spread foundations.} Rothermel derives a mathematical model for predicting rate of spread and
intensity in wildland fuels and remains the basis for many operational fire behavior tools \cite{rothermel1972spread}.
It provides fast, parameter-driven estimates suitable for real-time systems but is limited by assumptions about
fuels and homogeneity.

\textbf{Operational fire growth simulators.} Finney's FARSITE integrates multiple fire behavior components
(surface, crown, spotting, fuel moisture) to simulate two-dimensional fire growth over landscapes and has been
widely used for planning and decision support \cite{finney1998farsite}. FARSITE is valuable in DT settings because
it can be driven by GIS layers, but runtime can be a bottleneck for rapid re-forecasting without approximations.

\textbf{Physics-based models.} FIRETEC represents combustion, heat transfer, and atmosphere--fire interactions
using physics-based process models, enabling study of complex phenomena (wind transients, heterogeneous fuels)
\cite{linn2002firetec}. This higher fidelity comes with significantly higher computational cost, often requiring HPC
resources.

\textbf{Coupled atmosphere--fire forecasting.} Coen et al.\ present WRF-Fire, integrating a wildfire behavior
module within the Weather Research and Forecasting (WRF) model to capture two-way coupling between fire and local
meteorology \cite{coen2013wrffire}. Such coupling improves realism for wind-driven events but increases compute and
data requirements.

\textbf{Relevance to our project.} For a streaming prototype, we will prioritize fast simulators or simplified
spread envelopes (Rothermel-style or FARSITE-like) for responsiveness, while designing interfaces so higher-fidelity
coupled models (e.g., WRF-Fire) can be swapped in later.

\subsection{Observation products: satellite active-fire detection and fire-weather indices}
A wildfire DT is only as useful as its observation layer. Satellite active-fire products provide broad coverage,
while fire-weather indices summarize atmospheric drivers of fire potential.

\textbf{MODIS active fire (Collection 6).} Giglio et al.\ describe the MODIS Collection 6 fire detection algorithm and
product updates, including improved false-alarm control and revised fire radiative power retrieval \cite{giglio2016modis}.
MODIS offers global coverage with moderate spatial resolution and long historical continuity.

\textbf{VIIRS 375m active fire.} Schroeder et al.\ introduce the 375m VIIRS active fire product, providing finer spatial
detail and improved perimeter coherence compared with coarser products \cite{schroeder2014viirs}.

\textbf{Fire Weather Index (FWI).} Van Wagner's Canadian FWI System operationalizes weather-driven fire danger into
components linked to fuel moisture and spread potential \cite{vanwagner1987fwi}. Although regional calibration is
needed, FWI-like indices are widely used as features in risk models and for contextualizing alerts.

\textbf{Relevance to our project.} We will ingest VIIRS/MODIS active fire detections (as event streams) and compute
weather-derived risk features (FWI or similar) for spatiotemporal risk maps and for initializing/validating simple
spread forecasts.

\subsection{Data assimilation and real-time model updating}
\textbf{Data assimilation} updates model state (or parameters) using incoming measurements. For wildfire DTs,
assimilation is crucial because ignition location, fuel moisture, and winds are uncertain and evolve quickly.

\textbf{Ensemble Kalman filtering (EnKF) for wildfire.} Mandel et al.\ present a wildfire model with data assimilation
and demonstrate using ensemble-based methods to assimilate observations into running simulations \cite{mandel2008wildland}.

\textbf{Reduced-cost EnKF strategies.} Rochoux et al.\ propose approaches for computationally feasible wildfire spread
prediction with ensemble-based assimilation and uncertainty quantification \cite{rochoux2014enkf}.

\textbf{Perimeter assimilation with operational tools.} Yoo and Song develop a rapid method that assimilates wildfire
perimeters into a spread predictor using EnKF and accelerate updates via polyline simplification \cite{yoo2023rapid}.

\textbf{Relevance to our project.} In a course-scale prototype, we can implement a lightweight updating step
(e.g., perimeter correction or parameter nudging) compatible with a fast spread model, while logging uncertainty
proxies (ensemble spread or confidence scores) for downstream alerting.

\subsection{Cyberinfrastructure and decision-support platforms}
This category shows end-to-end systems integrating data, models, and user-facing products (maps, workflows,
ensembles) for operational use.

\textbf{Firemap (WIFIRE platform).} Crawl et al.\ present Firemap, a web platform that integrates real-time and
archived geospatial data with scalable processing to support simulation, visualization, and repeatable workflows,
including ensemble execution to represent uncertainty \cite{crawl2017firemap}.

\textbf{WIFIRE experience.} Altintas discusses cyberinfrastructure lessons from WIFIRE (data integration, workflow
automation, and operational deployment) as an example of translational impact \cite{altintas2021wifire}.

\textbf{Relevance to our project.} These systems validate the need for reproducible workflows and tight integration
between ingestion, analytics, modeling, and visualization. Our design goal is a smaller, open-source, course-feasible
subset: streaming ingestion + scalable geospatial analytics + simple forecasting + interactive dashboard.

\subsection{Big-data streaming and geospatial processing frameworks}
This category covers scalable engines and libraries that make real-time spatiotemporal analytics practical.

\textbf{Unified big-data processing.} Apache Spark provides a unified engine for large-scale batch and streaming
processing with a rich ecosystem \cite{zaharia2016spark}.

\textbf{Streaming ingestion backbone.} Kafka provides a durable distributed log for high-throughput event ingestion
and stream processing pipelines \cite{kreps2011kafka}.

\textbf{Geospatial analytics at scale.} GeoSpark/Sedona extends Spark with spatial data types, indexing, and joins
\cite{yu2019geospark}. ST-Hadoop extends MapReduce systems with spatiotemporal indexing and query support
\cite{alarabi2018sthadoop}.

\textbf{Relevance to our project.} Our pipeline will implement (i) streaming ingestion (Kafka-like topics),
(ii) processing (Spark Structured Streaming / micro-batches), and (iii) spatial joins (Sedona) to fuse active-fire
detections with contextual layers (weather grids, land cover, critical infrastructure proximity).

\section{Synthesis: gaps and how the proposed project fits}
Existing work provides strong components (satellite fire detections, spread models, assimilation methods, and
end-to-end platforms), but \textbf{small-team deployability} remains difficult due to integration complexity and the
operational burden of maintaining data pipelines, geospatial indexing, and continuous updating.

Our proposed project focuses on a \textbf{reproducible, lightweight architecture}:
\begin{itemize}
  \item \textbf{Streaming ingestion} of active-fire detections + weather snapshots + optional sensor data,
  \item \textbf{Scalable spatiotemporal joins} to compute risk layers and feature stores,
  \item \textbf{Simple, fast forecasting} (spread envelope / heuristic growth) with optional nudging,
  \item \textbf{Interactive dashboard + alerts} for decision support.
\end{itemize}
This design targets the practical gap between research prototypes and operational cyberinfrastructure.

\section{Summary table of categories and related work}
Table~\ref{tab:summary} highlights how each reference supports one or more categories relevant to the proposed
wildfire digital twin.

\begin{table}[h!]
\small
\centering
\renewcommand{\arraystretch}{1.15}
\begin{tabularx}{\textwidth}{lccccccX}
\toprule
\textbf{Ref} & \textbf{DT/DDDAS} & \textbf{Fire} & \textbf{Obs/RS} & \textbf{Assim.} & \textbf{Cyber} & \textbf{Big-Geo} & \textbf{Key idea (very short)} \\
\midrule
\cite{grieves2014digitaltwin} & \checkmark & & & & & & Digital twin concept baseline \\
\cite{darema2004dddas} & \checkmark & & & & & & DDDAS: online data into running sims \\
\cite{rothermel1972spread} & & \checkmark & & & & & Fast surface-spread rate model \\
\cite{finney1998farsite} & & \checkmark & & & & & Operational 2D fire growth simulator \\
\cite{linn2002firetec} & & \checkmark & & & & & Physics-based fire/atmosphere processes \\
\cite{coen2013wrffire} & & \checkmark & & & & & Coupled weather--fire forecasting \\
\cite{giglio2016modis} & & & \checkmark & & & & Active fire detection product \\
\cite{schroeder2014viirs} & & & \checkmark & & & & Finer-resolution fire detections \\
\cite{vanwagner1987fwi} & & & \checkmark & & & & Weather-based fire danger indices \\
\cite{mandel2008wildland} & & & & \checkmark & & & EnKF-style assimilation into wildfire model \\
\cite{rochoux2014enkf} & & & & \checkmark & & & Reduced-cost EnKF for spread prediction \\
\cite{yoo2023rapid} & & \checkmark & & \checkmark & & & Perimeter assimilation with acceleration \\
\cite{crawl2017firemap} & & & & & \checkmark & & End-to-end platform + workflows \\
\cite{altintas2021wifire} & & & & & \checkmark & & Cyberinfrastructure for response \\
\cite{halem2024wdt} & \checkmark & & & & \checkmark & & Wildfire DT vision and challenges \\
\cite{huang2024wdtprogress} & \checkmark & & \checkmark & & & & Wildfire DT progress review \\
\cite{zaharia2016spark} & & & & & & \checkmark & Unified engine for big data + streaming \\
\cite{kreps2011kafka} & & & & & & \checkmark & Distributed log + stream ingestion \\
\cite{yu2019geospark} & & & & & & \checkmark & Spatial joins/indexing on Spark \\
\cite{alarabi2018sthadoop} & & & & & & \checkmark & Spatio-temporal indexing on Hadoop \\
\bottomrule
\end{tabularx}
\caption{Mapping of related work to categories relevant to the proposed wildfire digital twin.}
\label{tab:summary}
\end{table}

\section*{Author Contributions}
\textbf{Prince Savsaviya:} Led problem formulation and system architecture; surveyed big-data streaming and geospatial
analytics literature (Spark, Kafka, Sedona/ST-Hadoop); integrated categories into a coherent narrative; edited final document.

\textbf{Teammate-1:} Surveyed wildfire spread modeling literature (Rothermel, FARSITE, WRF-Fire, FIRETEC) and summarized
model assumptions, strengths, and limitations.

\textbf{Ali:} Surveyed remote sensing and observation products (MODIS/VIIRS active fire, fire weather indices) and
documented how these data sources can be ingested and validated in real time.

\textbf{Teammate-3:} Surveyed data assimilation and real-time updating methods (EnKF variants, perimeter assimilation) and
wrote the synthesis of gaps/opportunities for a deployable DT prototype.

\bibliographystyle{ACM-Reference-Format}
\bibliography{refs}

\end{document}